\section{Methods}
\label{sec:methods}

% Mirrors docs/Formalism.md in the repository. If one changes, update the
% other.

\subsection{Problem statement}

Given a state $S$ partitioned into atomic units
$U = \{u_1, \ldots, u_M\}$ (census blocks, block groups, or VTDs) and
a target district count $N$, DualBalance Districting produces a
deterministic assignment $\pi : U \to \{1, \ldots, N\}$ such that
each district $D_i = \pi^{-1}(i)$ is contiguous, non-empty, and as
close as the geometry allows to representing both $1/N$ of the
state's people and $1/N$ of its geography. Let
$P = \sum_u \mathrm{pop}(u)$ and $A = \sum_u \mathrm{area}(u)$ with
per-district targets $P^* = P/N$ and $A^* = A/N$.

The full pipeline runs in three stages: a radial-seed stage,
\emph{PRISM} (\S\S\ref{sec:methods-seeds}--\ref{sec:methods-repair}),
which provides the deterministic starting configuration; a
VTD-scale tightening stage (\S\ref{sec:methods-optimize}) which
drives $\max_i |\delta_i|/P^{*}$ toward the user-supplied tolerance
$\tau$ via a hybrid greedy local search with bounded-chain escape;
and a block-scale refinement stage (\S\ref{sec:methods-block})
which re-initializes at finer granularity to recover area balance
within the achieved pop-deviation envelope. Each stage is a pure
function of its inputs.

\subsection{Radial seed placement}
\label{sec:methods-seeds}

Compute the population-weighted centroid
\[
  c = (c_x, c_y) =
  \left(
    \frac{\sum_u p_u\,x_u}{\sum_u p_u},\;
    \frac{\sum_u p_u\,y_u}{\sum_u p_u}
  \right),
\]
where $x_u, y_u$ are the centroid coordinates of unit $u$ in an
equal-area projection and $p_u$ its population. Let $\mathrm{diag}$
denote the bounding-box diagonal of the units and set the seed
radius $r = 0.001 \cdot \mathrm{diag}$. For $d = 0, 1, \ldots, N-1$
place seed $s_d$ at
\[
  s_d = \bigl(c_x + r\cos(2\pi d/N),\, c_y + r\sin(2\pi d/N)\bigr).
\]
Seed 0 sits due east of the centroid; seeds advance counter-clockwise
by equal angular steps. The choice $r = 0.001 \cdot \mathrm{diag}$ is
small enough that the Voronoi cells generated by the seeds degenerate
to near-perfect radial slices through $c$, but large enough to keep
the seed positions numerically distinct.

The radial configuration is what carries the dual-balance property:
each slice spans both the dense (near-$c$) and sparse (boundary-side)
territory of the state, so the population it inherits is bounded by
the cap (\S\ref{sec:methods-assign}) while the area it inherits is
driven toward $A^*$ by the slicing geometry.

\subsection{Capacitated first-fit assignment}
\label{sec:methods-assign}

Let $d(u, i) = \|x_u - s_i\|$ be the Euclidean distance from unit $u$
to seed $i$. Initialize remaining capacities $\rho_i = P^*$ for
$i = 1, \ldots, N$. Sort all $(u, i)$ pairs by ascending normalized
distance $d(u, i) / \mathrm{diag}$ and walk the sorted list in order:
\[
  \begin{array}{l}
    \text{for each } (u, i)\text{ in ascending } d(u,i)/\mathrm{diag}: \\
    \quad \text{if $u$ already assigned: continue} \\
    \quad \text{if $\rho_i \geq p_u$: assign $u$ to $i$;\;\; }\rho_i \leftarrow \rho_i - p_u \\
    \quad \text{else: skip}
  \end{array}
\]
Population balance is enforced as a hard cap: no district receives
more than $P^*$. Any unit not placed by the end of the walk (a rare
integer-rounding edge case) is assigned to the district with the
largest remaining capacity; \texttt{argmax} resolves ties to the
smallest district id.

Ties in normalized distance break by ascending
$(\mathrm{unit\_id},\,\mathrm{district\_id})$. There is no Lloyd
recentering, no iteration count, no convergence test: the radial
seeds do not drift, so a single assignment pass suffices.

\subsection{Contiguity repair}
\label{sec:methods-repair}

After \S\ref{sec:methods-assign} every unit belongs to exactly one
district, but a district may consist of more than one connected
component (rare on convex states; more common on those with
peninsulas, islands, or rural enclaves). For each such district,
the largest connected component by total population is retained;
units in the smaller components are reassigned to neighboring
districts one at a time, in the lowest-cost direction. Cost is
\[
  c(u, j) = \frac{d(x_u, s_j)}{\mathrm{diag}}
          + \frac{|P(D_j) + p_u - P^*|}{P^*}
          + \frac{|A(D_j) + a_u - A^*|}{A^*},
\]
combining a normalized distance term with normalized population and
area deviation penalties for the receiving district. Ties break in
cascade $(c, \mathrm{pop\_pen}, \mathrm{area\_pen}, \mathrm{dist},
\mathrm{district\_id})$ ascending. The repair sweep iterates until no
district has more than one connected component, capped at ten sweeps;
in practice it converges in zero or one sweep on real census
geometries.

\subsection{Scoring}

Define per-district relative deviations
$\mathrm{pop\_dev}_i = |P(D_i) - P^*|/P^*$ and
$\mathrm{area\_dev}_i = |A(D_i) - A^*|/A^*$ and let
$\overline{\mathrm{pop\_dev}}$, $\overline{\mathrm{area\_dev}}$
denote their means over $i = 1, \ldots, N$. The DualBalance Score is
\begin{equation}
  \mathrm{DBS}
    = \frac{1}{1
       + \tfrac{1}{2}\,\overline{\mathrm{pop\_dev}}
       + \tfrac{1}{2}\,\overline{\mathrm{area\_dev}}}.
  \label{eq:dbs}
\end{equation}
The $0.5/0.5$ weighting makes the error a convex combination of the
two mean deviations: each district is judged on representing roughly
$1/N$ of the people \emph{and} roughly $1/N$ of the state's geography.
The score reaches $1.0$ for a perfectly balanced plan and approaches
$0$ as deviations grow without bound. Secondary metrics
(Polsby-Popper compactness~\cite{polsbypopper1991} and
Reock~\cite{reock1961}) are computed alongside but not optimized
against; radial slices have lower compactness than blob-Voronoi or
hand-drawn districts by construction, and this is a deliberate trade
in service of the dual-balance objective.

PRISM does not directly minimize~\eqref{eq:dbs}; it minimizes
population-capacitated geographic assignment cost under radial
seeding. On the Minnesota PoC this still beats the enacted plan on
DBS ($0.6472$ vs $0.6390$) despite the lower compactness.

\subsection{Optional: \texorpdfstring{L\textsuperscript{1}}{L1} pop-tightening with DBS hill-climb}
\label{sec:methods-tighten}
\label{sec:methods-optimize}

The radial pipeline produces per-district pop deviation in the
5--15\% range on real census geometry, well above the
\textasciitilde{}0.5\% threshold required by \textit{Reynolds v.~Sims}
for U.S.~congressional districts~\cite{reynolds1964}. An optional
post-pass (\verb|--tighten-pop|) closes this gap via a deterministic
two-phase local search of boundary-unit moves. Let
$\delta_i = P(D_i) - P^{*}$ denote the signed population deviation of
district $D_i$, let $T$ denote the user-supplied tolerance (default
$0.005$), and call a move \emph{safe} if the source unit is not an
articulation point of its district's induced subgraph (so that
removing it preserves contiguity).

\paragraph{Phase 1: population tightening (hybrid L\textsuperscript{1} + max-norm).}
At each step, scan every boundary unit and every neighboring district
and compute both the L$^1$ change
\[
  \Delta_{L^1}(u, d_{\mathrm{dest}}) =
    |\delta_{d_{\mathrm{src}}} - p_u| - |\delta_{d_{\mathrm{src}}}|
  + |\delta_{d_{\mathrm{dest}}} + p_u| - |\delta_{d_{\mathrm{dest}}}|
\]
and the max-norm effect: a move \emph{strictly reduces} the global
maximum iff $|\delta_{d_{\mathrm{src}}} - p_u| < \max_i |\delta_i|$,
$|\delta_{d_{\mathrm{dest}}} + p_u| < \max_i |\delta_i|$, and at least
one of $\delta_{d_{\mathrm{src}}}, \delta_{d_{\mathrm{dest}}}$ is
itself at the current maximum. Each scanned candidate is labeled with
a priority key: priority $0$ if it strictly reduces the max,
priority $1$ if it only reduces L$^1$. Candidates are ranked
lexicographically by $(\text{priority}, \Delta_{L^1})$, and the safe
top-ranked candidate is accepted.

This hybrid is essential. Pure L$^1$-greedy leaves the worst district
untouched when it sits between two other over-target districts:
draining it requires L$^1$-neutral moves that the L$^1$ criterion
rejects. Pure max-norm gets stuck symmetrically: when two districts
are tied at the max, no single move strictly reduces the global
maximum (draining one leaves the other at the same value). The
hybrid takes any move that helps either norm, with max-reducing moves
preferred so the worst district is targeted whenever possible.

When neither 1-opt criterion finds an improving move but
$\max_i |\delta_i|/P^{*} > T$, invoke a bounded \emph{chain escape}:
search for an augmenting transport on the district-adjacency graph
of the form
\[
  u_0 : D_0 \to D_1,\; u_1 : D_1 \to D_2,\; \ldots,\; u_{k-1} : D_{k-1} \to D_k,
\]
of length $k = 2$ then $k = 3$, in which each $u_j$ is a boundary unit
between $D_j$ and $D_{j+1}$. The chain is accepted if it strictly
cuts the global maximum, or if it holds the maximum flat while
strictly reducing the L$^1$ sum (which is what happens when several
districts are tied at the max: each chain drains one of them, the
global maximum stays at the same value until the last tied district
is drained, but the L$^1$ sum drops monotonically). The unit-level
search inside each district triple sorts the arc lists $B_{i \to j}$
by population value; matching pairs $(u, v)$ are found via binary
search around a target population value rather than by Cartesian
enumeration, bounding the per-triple cost by
$O(|B_{i \to j}| \log |B_{j \to k}|)$. Application is in
\emph{reverse} order ($u_{k-1}$ first, then $u_{k-2}$, etc.) so the
intermediate districts never temporarily overload during the chain.
This is the deterministic analogue of an ejection chain. Phase 1
terminates only when neither 1-opt nor bounded-chain escape finds an
improving sequence.

The L$^1$ objective is essential to the radial geometry. Seed
placement puts the most over-target and most under-target districts
on opposite sides of the population centroid, so no single
adjacent-slice swap reduces the L$^\infty$ maximum, but many such
swaps reduce the sum. The canonical Reynolds-tightening literature
uses L$^\infty$ and bottoms out at \textasciitilde{}5\% on this
geometry; the hybrid pass with chain escape runs to completion in
\textasciitilde{}80 swaps on Minnesota, reducing
$\mathrm{pop\_dev\_max}$ from $0.1124$ to $0.0021$ while leaving
$\mathrm{area\_dev\_mean}$ essentially unchanged.

\paragraph{Phase 2: DBS hill-climb.}
Once Phase 1 has converged, at each step pick the safe boundary move
that maximally improves $\mathrm{DBS}$ (equation~\ref{eq:dbs}), subject
to $\max_i |\delta_i|/P^{*}$ not exceeding the value Phase 1 reached
(equivalently, never crossing back above $T$). Phase 2 runs until no
improving safe move exists. Whereas the core PRISM pass does not
directly minimize equation~\ref{eq:dbs}, this opt-in Phase 2 does, but
only on the post-Phase-1 plan and only within the Reynolds-compliant
feasible region.

\paragraph{Determinism and engineering.}
Both phases are pure functions of $(\mathrm{units},\,N,\,T)$. The only
source of dependency beyond the core pipeline's $(\mathrm{units},\,N)$
is the explicit tolerance $T$. The optimizer maintains a per-district
articulation-point cache via Tarjan's algorithm on CSR adjacency
arrays~\cite{tarjan1972}, reducing the per-candidate contiguity check
from $O(V+E)$ to $O(1)$; an incrementally tracked boundary-unit set
restricts each scan to units that have a different-district neighbor.
At VTD scale the speedup is modest; at block scale (tens of thousands
of units per district) it is the difference between hours and minutes.

The pass is off by default and gated by an explicit flag. Pure radial
remains the principled algorithm; the two-phase optimizer is an opt-in
concession to legal compliance, with Phase 2 providing the further
courtesy of directly hill-climbing the project's stated objective once
the Reynolds constraint is satisfied.

\subsection{Multi-resolution refinement: VTD to Census block}
\label{sec:methods-block}

The hybrid Phase 1 + chain escape closes most of the Reynolds gap
at VTD scale, but not all of it. Modern U.S. congressional plans
must clear the much tighter \emph{Karcher v.~Daggett}~\cite{karcher1983}
practical threshold of $\sim 0.05\%$ total deviation, and at VTD
granularity the residual gap is structural rather than
algorithmic. The smallest single-unit move has the size of the
smallest VTD's population. On real states the median VTD population
is on the order of $10^3$, while the Karcher budget on a $\sim 800$k
ideal district is $\sim 400$ people. Most candidate moves overshoot
the budget on at least one endpoint, and Phase 2 starves: the
running-max constraint admits few moves, and the algorithm settles
above Karcher with most of the residual area-balance gain
unrealized.

The fix is to refine at finer granularity. We re-run the optimizer
on Census tabulation blocks (median population $\sim 20$ rather
than $\sim 10^3$), initializing from the VTD-Karcher plan rather
than re-seeding from scratch. The pipeline is:

\begin{enumerate}
  \item Compute the VTD-level plan
        $\pi_{\mathrm{VTD}}: V_{\mathrm{VTD}} \to \{1, \ldots, N\}$
        via PRISM + Phase 1 + Phase 2 at VTD scale with tolerance $T$.
  \item Load the block-level units $V_{\mathrm{block}}$ and project
        both layers to the same equal-area CRS.
  \item For each block $u \in V_{\mathrm{block}}$, find the VTD
        polygon $v(u) \in V_{\mathrm{VTD}}$ containing its
        representative point, and set
        $\pi_0(u) := \pi_{\mathrm{VTD}}(v(u))$. Blocks whose
        representative point lies outside every VTD polygon (rare,
        only on coastal/water boundary cases) fall back to the
        nearest VTD by centroid distance. The spatial join is
        deterministic given a fixed tie-breaking order on
        $\mathrm{unit\_id}$.
  \item Run Phase 1 + Phase 2 again at block scale starting from
        $\pi_0$, with the same tolerance $T$.
\end{enumerate}

Because VTDs are unions of whole blocks, $\pi_0$ inherits the
VTD-level plan's pop totals and contiguity exactly: the block-level
plan is identical to the VTD-level plan re-described at finer
granularity. Phase 1 at block scale therefore has nothing to do on
states where the VTD-level pass reached $T$. Phase 2 has substantial
freedom: block populations are two orders of magnitude smaller than
VTD populations, so the running-max envelope around $T$ admits many
candidate moves and the DBS hill-climb runs until saturation.

Because PRISM-style seeding at block scale starts from a much worse
initial configuration than the VTD seed (the population-weighted
centroid is the same but the much higher density of degenerate
articulation-point boundaries traps the optimizer), the refinement
strategy is the path that works in practice: VTD as a coarse
solver, blocks for the precise area-balance recovery.

\subsection{Engineering: per-district articulation-point cache}
\label{sec:methods-contig}

The hot loop of the optimizer asks repeatedly whether a candidate
boundary unit $u$ can be safely removed from its district $d$ -- in
other words, whether $u$ is an articulation point of the subgraph
$G[V_d]$. A naive implementation runs $\mathcal{O}(|V_d| + |E_d|)$
per query via \verb|networkx.is_connected| on the post-removal
subgraph, which dominates runtime at block scale.

We maintain a per-district cache. For each district $d$, the set
$A_d \subseteq V_d$ of articulation points is precomputed by Tarjan's
algorithm~\cite{tarjan1972} on CSR (compressed sparse row) integer
adjacency arrays. The per-query \verb|can_remove(u)| reduces to the
set membership $u \notin A_{\pi(u)}$, returned in $\mathcal{O}(1)$
amortized. After each accepted move $u: d_{\mathrm{src}} \to d_{\mathrm{dest}}$
the two affected articulation sets $A_{d_{\mathrm{src}}}$,
$A_{d_{\mathrm{dest}}}$ are recomputed; districts not touched by the
move keep their cached set unchanged.

Each Tarjan recompute is $\mathcal{O}(|V_d| + |E_d|)$ rather than
the global $\mathcal{O}(|V| + |E|)$ of the naive check
($|V_d| \approx |V|/N$ for $N$ districts). The recompute itself is
\verb|@njit|-compiled via numba when available, operating on
preallocated workspace buffers (a single integer DFS stack, two
boolean masks, and three integer arrays sized to $|V|$). On a
typical block-scale district subgraph of $\sim 44{,}000$ nodes the
JIT'd recompute completes in $\sim 20$~ms, compared to $\sim 560$~ms
for the same query through \verb|networkx.articulation_points|: a
$28\times$ speedup. Across a full block-scale optimizer run
(thousands of accepted moves and tens of thousands of candidates per
move), this is the difference between a 20-hour run and a 30-minute
run.

A boundary-unit set $\partial = \{u : \exists v \in N(u), \pi(v) \neq \pi(u)\}$
is also maintained incrementally and updated after each move (only
$u$ and its graph neighbors need to be re-checked). The inner scan
iterates $\partial$ rather than $V$, dropping per-pass work from
$|V|$ to $O(|\partial|)$, which on real states is roughly the
square root of $|V|$.

Neither caching layer changes the algorithm's output. They preserve
determinism exactly and make the block-scale pipeline tractable on
commodity hardware.

\subsection{Determinism and tie-breaking}

The pipeline is deterministic at every step. The only sources of
ambiguity (equidistant seed-to-unit pairs in
\S\ref{sec:methods-assign}, equal-capacity fallbacks for unplaced
units, and equal-cost candidates in \S\ref{sec:methods-repair}) all
resolve to a fixed cascade with no remaining ambiguity:

\begin{enumerate}
  \item In assignment, ties on normalized distance break by ascending
        $(\mathrm{unit\_id}, \mathrm{district\_id})$.
  \item In the fallback for unplaced units, ties on remaining capacity
        break to the smallest district id.
  \item In contiguity repair, ties on cost break by ascending
        $(\mathrm{pop\_pen}, \mathrm{area\_pen}, \mathrm{distance},
        \mathrm{district\_id})$.
\end{enumerate}

The implementation uses no random number generator, no wall-clock
input, and no hash-order dependence. Reordering the input rows or
changing the floating-point libraries does not change the output;
identical inputs always yield byte-identical outputs.

\subsection{Out-of-scope inputs}

The generator reads only geography and population. Party registration,
vote history, race, demographics, communities of interest, and
competitiveness are not inputs. The scoring harness may \emph{report}
partisan or demographic diagnostics on the resulting map but does not
feed them back into the generator.
